\documentclass[11pt,a4paper]{article}

% Pakete
\usepackage[utf8]{inputenc}
\usepackage[T1]{fontenc}
\usepackage[ngerman]{babel}
\usepackage[left=2cm,right=2cm,top=2cm,bottom=2cm]{geometry}
\usepackage{helvet}
\renewcommand{\familydefault}{\sfdefault}
\usepackage{amsmath,amssymb}
\usepackage{graphicx}
\usepackage{tikz}
\usetikzlibrary{arrows.meta,positioning}
\usepackage{tcolorbox}
\tcbuselibrary{skins,breakable}
\usepackage{enumitem}
\usepackage{tabularx}
\usepackage{colortbl}
\usepackage{siunitx}
\sisetup{locale=DE, per-mode=fraction}
\usepackage{qrcode}

% Fachfarbe Physik
\definecolor{physik}{HTML}{2c5aa0}
\definecolor{physiklight}{HTML}{e3f2fd}

% Box-Stile
\tcbset{
    merkbox/.style={
        colback=physiklight,
        colframe=physik,
        fonttitle=\bfseries,
        title=#1,
        boxrule=1pt,
        arc=2pt,
        left=5pt,
        right=5pt,
        top=5pt,
        bottom=5pt
    }
}

% Kopfzeile
\usepackage{fancyhdr}
\pagestyle{fancy}
\fancyhf{}
\lhead{\small Physik 12 GK}
\chead{\small Photoeffekt -- Erkenntnisgewinnung}
\rhead{\small Name: \underline{\hspace{4cm}}}
\renewcommand{\headrulewidth}{0.4pt}

% Keine Einrückung
\setlength{\parindent}{0pt}
\setlength{\parskip}{6pt}

\begin{document}

% Titel mit QR-Codes
\begin{minipage}[t]{0.65\textwidth}
\vspace{0pt}
{\Large\textbf{Der Photoeffekt -- Einsteinsche Gerade}}\\[0.3em]
{\normalsize Stunde 2: Vom Experiment zur Theorie}
\end{minipage}
\hfill
\begin{minipage}[t]{0.32\textwidth}
\vspace{0pt}
\raggedleft
\begin{tabular}{@{}c@{\hspace{8pt}}c@{}}
\qrcode[height=1.4cm]{https://phet.colorado.edu/de/simulations/photoelectric} &
\qrcode[height=1.4cm]{https://example.com/LP-02} \\
{\scriptsize PhET-Sim} & {\scriptsize Lernpfad}
\end{tabular}
\end{minipage}

\vspace{0.5em}

% ============================================================
% KASTEN 1: AUSGANGSSITUATION
% ============================================================

\begin{tcolorbox}[merkbox={Kasten 1: Ausgangssituation}]
\textbf{Problem:} Das Wellenmodell erklärt den Photoeffekt nicht.\\
\textbf{Hypothese (Einstein 1905):} Licht besteht aus Energiepaketen -- \textbf{Photonen} mit $E = h \cdot f$.

\vspace{0.3em}
\textbf{Frage:} Wie hängt die Energie der ausgelösten Elektronen von der Lichtfrequenz ab?
\end{tcolorbox}

\vspace{0.3em}

% ============================================================
% KASTEN 2: GEGENFELDMETHODE
% ============================================================

\begin{tcolorbox}[merkbox={Kasten 2: Gegenfeldmethode}]
\textbf{Prinzip:} Ein elektrisches Feld bremst die Elektronen ab.

Die \textbf{Grenzspannung} $U_G$ stoppt gerade die schnellsten Elektronen:
\hfill $\boxed{E_{\text{kin,max}} = e \cdot U_G}$ \hfill\mbox{}

\vspace{0.3em}
\textbf{Praktisch:} Bei $U_G$ in Volt ist $E_{\text{kin}}$ numerisch gleich in \textbf{Elektronenvolt (eV)}.
\end{tcolorbox}

\vspace{0.5em}

% ============================================================
% AUFGABE 1: MESSREIHE
% ============================================================

\textbf{Aufgabe 1: Messreihe aufnehmen} \hfill \textbf{(10 BE)}

Öffne die PhET-Simulation (QR-Code). Wähle \textbf{Natrium} und miss die Grenzspannung für verschiedene Wellenlängen.

\vspace{0.3em}

\begin{tabular}{|c|c|c|c|}
\hline
\textbf{Wellenlänge $\lambda$ (nm)} & \textbf{Frequenz $f$ ($\SI{e14}{Hz}$)} & \textbf{Grenzspannung $U_G$ (V)} & \textbf{$E_{\text{kin}}$ (eV)} \\
\hline
400 (violett) & 7,5 & & \\[0.8em]
\hline
450 (blau) & 6,7 & & \\[0.8em]
\hline
500 (grün) & 6,0 & & \\[0.8em]
\hline
550 (gelb) & 5,5 & & \\[0.8em]
\hline
600 (orange) & 5,0 & & \\[0.8em]
\hline
\end{tabular}

\textit{Hinweis: Bei manchen Wellenlängen gibt es keinen Photoeffekt -- trage dann ,,---'' ein.}

\vspace{0.8em}

% ============================================================
% KASTEN: INTERPRETATION DER GERADEN
% ============================================================

\begin{tcolorbox}[merkbox={Vorüberlegung: Was verrät uns eine Gerade?}]
Wenn $E_{\text{kin}}$ linear von $f$ abhängt, gilt: \quad $E_{\text{kin}} = m \cdot f + n$

\vspace{0.3em}
\begin{tabular}{@{}ll@{}}
\textbf{Steigung} $m$ & = Änderung von $E_{\text{kin}}$ pro Frequenzeinheit \\
\textbf{y-Achsenabschnitt} $n$ & = Wert von $E_{\text{kin}}$ bei $f = 0$ (Extrapolation!) \\
\textbf{x-Achsenabschnitt} $f_G$ & = Frequenz, bei der $E_{\text{kin}} = 0$ wird \\
\end{tabular}

\vspace{0.3em}
\textit{Überlege: Was bedeuten diese Größen physikalisch beim Photoeffekt?}
\end{tcolorbox}

\newpage

% ============================================================
% AUFGABE 2: DIAGRAMM (große Darstellung mit 5mm-Gitter)
% ============================================================

\textbf{Aufgabe 2: Grafische Darstellung} \hfill \textbf{(8 BE)}

Trage deine Messwerte ein. Zeichne eine Ausgleichsgerade und \textbf{extrapoliere} sie bis zur y-Achse.

\textit{Beschrifte die Achsen selbst: x-Achse bis 10, y-Achse von $-3$ bis $+2$.}

\vspace{0.3em}

\begin{center}
\begin{tikzpicture}[scale=1.0]
    % 5mm Gitter (bei scale=1.0 entspricht 1cm = 1 Einheit)
    % Gitter von x: 0-10, y: -3 bis +2
    \draw[gray!40, very thin, step=0.5] (0,-3) grid (10,2);

    % Achsen (ohne Beschriftung - Schüler machen das selbst)
    \draw[->, thick] (0,-3) -- (0,2.3);
    \draw[->, thick] (0,0) -- (10.3,0);

    % Achsenbeschriftung nur Variablen
    \node[right] at (10.3,0) {$f$};
    \node[above] at (0,2.3) {$E_{\text{kin}}$};

    % Nullpunkt markieren
    \node[below left] at (0,0) {\small 0};
\end{tikzpicture}
\end{center}

\vspace{0.8em}

% ============================================================
% AUFGABE 3: AUSWERTUNG
% ============================================================

\textbf{Aufgabe 3: Auswertung der Einsteinschen Geraden} \hfill \textbf{(12 BE)}

\begin{enumerate}[label=\alph*)]
\item \textbf{Steigung bestimmen:} Lies zwei Punkte auf deiner Geraden ab und berechne die Steigung. \hfill (4 BE)

\vspace{0.3em}
Punkt 1: $f_1 = $ \underline{\hspace{1.5cm}} $\cdot 10^{14}$ Hz \quad $E_1 = $ \underline{\hspace{1.5cm}} eV

Punkt 2: $f_2 = $ \underline{\hspace{1.5cm}} $\cdot 10^{14}$ Hz \quad $E_2 = $ \underline{\hspace{1.5cm}} eV

\vspace{0.5em}

Steigung: $m = \dfrac{\Delta E}{\Delta f} = \dfrac{E_2 - E_1}{f_2 - f_1} = $ \underline{\hspace{3cm}} $= $ \underline{\hspace{2cm}} \si{eV.s}

\vspace{0.3em}
\textit{Dies ist das} \textbf{Plancksche Wirkungsquantum} $h$ (Literaturwert: $h = \SI{4,14e-15}{eV.s}$).

\vspace{0.8em}

\item \textbf{x-Achsenabschnitt:} Wo schneidet deine Gerade die x-Achse? \hfill (2 BE)

$f_G = $ \underline{\hspace{2cm}} $\cdot 10^{14}$ Hz

\vspace{0.3em}
Physikalische Bedeutung: \underline{\hspace{9cm}}

\vspace{0.8em}

\item \textbf{Austrittsarbeit:} Stelle die Photogleichung $E_{\text{kin}} = h \cdot f - W_A$ nach $W_A$ um \hfill (2 BE)

und setze $E_{\text{kin}} = 0$ bei $f = f_G$ ein:

\vspace{0.8cm}

$W_A = $ \underline{\hspace{3cm}} $\approx$ \underline{\hspace{2cm}} eV

\vspace{0.8em}

\item \textbf{y-Achsenabschnitt:} Lies ab, wo deine Gerade die y-Achse schneidet. \hfill (2 BE)

$n = $ \underline{\hspace{2cm}} eV

\vspace{0.3em}
Was bedeutet dieser negative Wert physikalisch? \underline{\hspace{6cm}}

\vspace{0.8em}

\item \textbf{Geradengleichung:} Formuliere die vollständige Gleichung. \hfill (2 BE)

$E_{\text{kin}} = $ \underline{\hspace{1.5cm}} $\cdot f - $ \underline{\hspace{1.5cm}}

\end{enumerate}

\vspace{0.5em}

% ============================================================
% KASTEN 3: ZUSAMMENFASSUNG
% ============================================================

\begin{tcolorbox}[merkbox={Kasten 3: Die Photogleichung (Ergebnis)}]
\textbf{Energiebilanz beim Photoeffekt:}

\vspace{0.3em}
\hfill $\boxed{E_{\text{kin,max}} = h \cdot f - W_A}$ \hfill mit \hfill $h = \SI{4,14e-15}{eV.s}$ \hfill\mbox{}

\vspace{0.5em}
\begin{tabular}{ll}
$W_A$ & = \textbf{Austrittsarbeit} (materialabhängig) \\
$f_G = W_A / h$ & = \textbf{Grenzfrequenz} (unterhalb: kein Photoeffekt) \\
\end{tabular}
\end{tcolorbox}

\vspace{0.3em}

% ============================================================
% AUFGABE 4: REFLEXION
% ============================================================

\textbf{Aufgabe 4: Reflexion} \hfill \textbf{(4 BE)}

\begin{enumerate}[label=\alph*)]
\item Warum konnte das Plancksche Wirkungsquantum $h$ aus dem Experiment bestimmt werden? \hfill (2 BE)

\vspace{1cm}

\item Nenne zwei mögliche Fehlerquellen bei deiner h-Bestimmung. \hfill (2 BE)

\vspace{1cm}

\end{enumerate}

\begin{center}
\small\textit{Gesamtpunktzahl: 34 BE}
\end{center}

\end{document}
